\chapter{La Cascade Hormonale : Le Corps qui Hurle}

\textit{Nous pensons souvent que nos paroles tues s'évaporent dans l'air. C'est une erreur de perspective. Toute parole retenue par peur est une pierre jetée dans le puits de notre physiologie. Elle ne disparaît pas ; elle coule au fond et trouble l'eau de notre santé.}

\section{L'Inhibition de l'Action : Le Piège de Laborit}
Henri Laborit l'avait pressenti : ce n'est pas l'agression qui nous rend malades, c'est l'impossibilité d'y répondre. Pour l'homme ou la femme incapable de s'affirmer, la vie quotidienne est une suite de "petites agressions" sans réponse. Le cerveau ordonne le combat, l'adrénaline inonde le sang, le cœur s'emballe... et rien ne se passe. La parole reste coincée entre les dents. Cette énergie, faute de pouvoir sortir sous forme de mots, se retourne contre nos propres organes.

\section{Le Cortisol : L'Acide de la Soumission}
Là où l'adrénaline est le feu de l'instant, le cortisol est la braise qui consume. Dans le silence de l'inhibition, le cortisol ne redescend jamais. Il devient un poison lent. Il circule dans les veines, fragilise l'immunité, et sature nos récepteurs neuronaux. 

\begin{NoteClinique}
\textbf{Le cas de Sophie, 35 ans.} \\
Sophie ne sait pas dire "non". Ni à sa mère, ni à son patron, ni à ses amies qui abusent de sa gentillesse. Physiologiquement, Sophie vit dans un état de "sympathicotonie" permanente. Ses nuits sont peuplées d'insomnies car son corps est en alerte, attendant une décharge qui ne vient jamais. Elle a développé un psoriasis sévère sur les mains. Ce n'est pas "dans sa tête" ; c'est son système nerveux qui, à force d'être sous haute tension, finit par "griller" ses propres circuits cutanés.
\end{NoteClinique}

\section{Le Ventre, Théâtre de l'Inhibition}
Le nerf vague, ce lien mystérieux entre notre cerveau et nos viscères, est le premier impacté par le manque d'affirmation. Lorsque nous nous taisons par peur, nous envoyons un signal de danger à notre système digestif. Le "nœud à l'estomac" n'est pas une image littéraire, c'est une réalité neurologique. 
